% NOETHER PAPER 2 — KOREAN TRANSLATION-PRODUCER DRAFT — UNCHECKED
% Unit: T03-U10; current-authority whole lines 599--603
% Source slice: 1,125 LF UTF-8 bytes; SHA-256 BF24F4A6AC52C0E4594E4618C22668C78D57AF6D2E94FCCFFCB307EFF6A02E5D
% Pointer: NOETH-DE-AUTH-v004-20260804; SHA-256 A1C62FDACAA34DFC1B806DC18258F2E732539F7AE9D85AA4BD9E1067B8749D9F
% This is producer translation only: no source/Korean/formula review, compilation, or rendering.

\emph{형식열}\srcfn{**)}{\emph{Clebsch}, Über eine Fundamentalaufgabe der Invariantentheorie (Abh. der Gött. Ges. d. Wiss. Bd. XVII), § 17 및 § 18 끝 참조. \glqq{}형식열\grqq{}은 더 높은 형식에 따른 배열 원리에 의해 그곳에서 도입한 \glqq{}본래적으로 환원된 동치계\grqq{}와 구별된다.}\setcounter{footnote}{9}을 하나의 시작형식과 이 시작형식의 자기 수축으로 생긴 모든 형식을 합한 것으로 정의하고, 그 시작형식으로 형식열을 나타낸다. 여기서 자기 수축이 더 많은 형식을 \glqq{}더 높은 형식\grqq{}이라 하며, 서로 동등한 수축 $s_\tau$와 $t_\sigma$ 가운데 하나를 더 높은 것으로 정규화해야 한다. 형식열의 구성 법칙에서 다음이 따른다.

1) 형식열의 각 형식은 시작형식의 기호들에 관해 선형이다.

2) 서로 반변인 기호열을 각각 두 개씩 가진 시작형식에서는 형식열이 유일하게 정해진다. 더 많은 기호열을 가진 시작형식에서는 항등식 정리로 연결되는 수축들 가운데 특정 수축들을 우선시함으로써 형식열을 유일하게 정규화해야 한다.

